% GENERUOJAMA is rezultatai/darbiniai/nematytos_test.csv
% Ranka NELIESTI - paleisti: python -m src.eksperimentai.nematytos
\begingroup
\footnotesize
\setlength{\tabcolsep}{5pt}
\begin{tabularx}{\textwidth}{@{}>{\raggedright\arraybackslash}X>{\raggedleft\arraybackslash}p{1.3cm}>{\centering\arraybackslash}p{2.1cm}>{\centering\arraybackslash}p{2.1cm}>{\centering\arraybackslash}p{2.1cm}@{}}
\toprule
\multirow{2}{*}{\textbf{Pašalinta klasė}} & \multirow{2}{*}{\textbf{$n$}} & \multicolumn{3}{c}{\textbf{Aptikta (pažymėta kaip ataka)}} \\
\cmidrule(l){3-5}
 & & \textbf{Prižiūrimas, klasės nematęs} & \textbf{Autokoderis} & \textbf{Prižiūrimas, klasę matęs} \\
\midrule
\texttt{DDOS-SLOWLORIS} & 3\,360 & 99,9~\% & 31,4~\% & 99,9~\% \\
\texttt{RECON-PORTSCAN} & 11\,518 & 47,2~\% & 12,8~\% & 48,7~\% \\
\texttt{DICTIONARYBRUTEFORCE} & 1\,878 & 24,4~\% & 25,3~\% & 45,7~\% \\
\bottomrule
\end{tabularx}
\vspace{2pt}
\raggedright\scriptsize Testavimo aibė. Klausiama tik to, ar eilutė pažymima kaip bet kuri ataka: pašalinus \texttt{DICTIONARYBRUTEFORCE} ištuštėja visa \texttt{BruteForce} kategorija, todėl macro-F1 su baziniu modeliu būtų nepalyginamas. Autokoderis nepermokomas, nes mokomas tik iš gerybinio srauto, tad visos atakų klasės jam ir taip nematytos. Paskutinis stulpelis rodo, kiek kainuoja klasės nematyti. Kiekvienam modeliui $\tau$ parinktas validacijos aibėje.
\endgroup
